% GENERATED FILE. Edit canonical lessons, then run npm run generate:books.
% canonical-source-hash: fnv1a64:d29c5e46986fccaa
% canonical-lessons: MW-C218-rel, MW-C218-savari, MW-C218-draivar, MW-C218-chhutno, MW-C218-pugavno

\chapter{Bus and Train}
\label{ch:mw-bus-and-train}
\hlchaptermodality{218}

\begin{chapteropening}
\textbf{By the end of this chapter:} \emph{I can say I will take the train and when the bus leaves.}

Complete the last lesson of chapter 218: I can say I will take the train and when the bus leaves.
\end{chapteropening}

\section[\texorpdfstring{\mw{रेल} (rel)}{रेल (rel)}]{\mw{रेल} (rel) — a train}

\label{lesson:MW-C218-rel}



\begin{quote}
Before the new one: say the Marwadi for to load, then the Marwadi for to halt.
\end{quote}

\subsection*{You'll want to know: \mw{रेल}}
\textbf{\mw{रेल}} — \emph{rel} — \textquotedblleft{}a train\textquotedblright{}.

\begin{quote}
\mw{हूं} \mw{काल} \mw{रेल} \mw{मांय} \mw{बैठसूं।}
\end{quote}

\begin{quote}
\emph{hū̃ kāl rel mā̃y baiṭhsū̃.}
\end{quote}

\begin{quote}
Tomorrow I'll take the train.
\end{quote}

A feminine noun.

\begin{cousinweb}
\textbf{\mw{रेल}} (\emph{rel}) is borrowed from English \emph{rail}; Marwadi uses it for the train itself.
\end{cousinweb}

\begin{grammarlens}[title={What you've built}]
One more word for saying what comes next.
\end{grammarlens}

\subsection*{Your turn}
\begin{itemize}
\raggedright
  \item \emph{Say it:} \emph{rel}
  \item \emph{Say it:} \emph{rel}, once more
  \item \emph{Say it:} \emph{hū̃ kāl rel mā̃y baiṭhsū̃.} — Tomorrow I'll take the train.
  \item \emph{Recall:} say the Marwadi for to load, then the Marwadi for to halt, then say \emph{rel} again
\end{itemize}

\begin{tcolorbox}[breakable,colback=teal!4,colframe=teal!35!black,title={Before you move on}]
What does \mw{रेल} mean? (\textquotedblleft{}A train\textquotedblright{}.) Say it once more.
\end{tcolorbox}
\section[\texorpdfstring{\mw{सवारी} (savārī)}{सवारी (savārī)}]{\mw{सवारी} (savārī) — a passenger; a ride}

\label{lesson:MW-C218-savari}



\begin{quote}
Before the new one: say the Marwadi for to halt, then the Marwadi for a train.
\end{quote}

\subsection*{You'll want to know: \mw{सवारी}}
\textbf{\mw{सवारी}} — \emph{savārī} — \textquotedblleft{}a passenger; a ride\textquotedblright{}.

\begin{quote}
\mw{सवारी} \mw{बस} \mw{मांय} \mw{बैठै} \mw{है।}
\end{quote}

\begin{quote}
\emph{savārī bas mā̃y baiṭhai hai.}
\end{quote}

\begin{quote}
The passenger sits in the bus.
\end{quote}

A feminine noun.

\begin{cousinweb}
\textbf{\mw{सवारी}} (\emph{savārī}) is usually traced to Persian \emph{savārī}, \textquotedblleft{}riding, a mount\textquotedblright{}; in Marwadi it has come to mean the one who rides, a passenger.
\end{cousinweb}

\begin{grammarlens}[title={What you've built}]
One more word for saying what comes next.
\end{grammarlens}

\subsection*{Your turn}
\begin{itemize}
\raggedright
  \item \emph{Say it:} \emph{savārī}
  \item \emph{Say it:} \emph{savārī}, once more
  \item \emph{Say it:} \emph{savārī bas mā̃y baiṭhai hai.} — The passenger sits in the bus.
  \item \emph{Recall:} say the Marwadi for to halt, then the Marwadi for a train, then say \emph{savārī} again
\end{itemize}

\begin{tcolorbox}[breakable,colback=teal!4,colframe=teal!35!black,title={Before you move on}]
What does \mw{सवारी} mean? (\textquotedblleft{}A passenger; a ride\textquotedblright{}.) Say it once more.
\end{tcolorbox}
\section[\texorpdfstring{\mw{ड्राइवर} (ḍrāivar)}{ड्राइवर (ḍrāivar)}]{\mw{ड्राइवर} (ḍrāivar) — a driver}

\label{lesson:MW-C218-draivar}



\begin{quote}
Before the new one: say the Marwadi for a train, then the Marwadi for a passenger.
\end{quote}

\subsection*{You'll want to know: \mw{ड्राइवर}}
\textbf{\mw{ड्राइवर}} — \emph{ḍrāivar} — \textquotedblleft{}a driver\textquotedblright{}.

\begin{quote}
\mw{ड्राइवर} \mw{दिनूगै} \mw{आवै} \mw{है।}
\end{quote}

\begin{quote}
\emph{ḍrāivar dinūgai āvai hai.}
\end{quote}

\begin{quote}
The driver comes early in the morning.
\end{quote}

A masculine noun.

\begin{cousinweb}
\textbf{\mw{ड्राइवर}} (\emph{ḍrāivar}) is borrowed from English \emph{driver}.
\end{cousinweb}

\begin{grammarlens}[title={What you've built}]
One more word for saying what comes next.
\end{grammarlens}

\subsection*{Your turn}
\begin{itemize}
\raggedright
  \item \emph{Say it:} \emph{ḍrāivar}
  \item \emph{Say it:} \emph{ḍrāivar}, once more
  \item \emph{Say it:} \emph{ḍrāivar dinūgai āvai hai.} — The driver comes early in the morning.
  \item \emph{Recall:} say the Marwadi for a train, then the Marwadi for a passenger, then say \emph{ḍrāivar} again
\end{itemize}

\begin{tcolorbox}[breakable,colback=teal!4,colframe=teal!35!black,title={Before you move on}]
What does \mw{ड्राइवर} mean? (\textquotedblleft{}A driver\textquotedblright{}.) Say it once more.
\end{tcolorbox}
\section[\texorpdfstring{\mw{छूटणो} (chhūṭṇo)}{छूटणो (chhūṭṇo)}]{\mw{छूटणो} (chhūṭṇo) — to depart (of a bus or a train); to get loose}

\label{lesson:MW-C218-chhutno}



\begin{quote}
Before the new one: say the Marwadi for a passenger, then the Marwadi for a driver.
\end{quote}

\subsection*{You'll want to know: \mw{छूटणो}}
\textbf{\mw{छूटणो}} — \emph{chhūṭṇo} — \textquotedblleft{}to depart (of a bus or a train); to get loose\textquotedblright{}.

\begin{quote}
\mw{बस} \mw{आथण} \mw{छूटै} \mw{है।}
\end{quote}

\begin{quote}
\emph{bas āthaṇ chhūṭai hai.}
\end{quote}

\begin{quote}
The bus leaves in the evening.
\end{quote}

Past (it …): \mw{वो} \mw{बस} \mw{छूटी}, \emph{vo bas chhūṭī}.

\begin{cousinweb}
\textbf{\mw{छूटणो}} (\emph{chhūṭṇo}) is inherited from a reconstructed \emph{*kṣuṭyate}, \textquotedblleft{}is released\textquotedblright{}; CDIAL lists the Marwari verb, \textquotedblleft{}to get loose\textquotedblright{}, under it (CDIAL 3707).
\end{cousinweb}

\begin{grammarlens}[title={What you've built}]
One more word for saying what comes next.
\end{grammarlens}

\subsection*{Your turn}
\begin{itemize}
\raggedright
  \item \emph{Say it:} \emph{chhūṭṇo}
  \item \emph{Say it:} \emph{chhūṭṇo}, once more
  \item \emph{Say it:} \emph{bas āthaṇ chhūṭai hai.} — The bus leaves in the evening.
  \item \emph{Recall:} say the Marwadi for a passenger, then the Marwadi for a driver, then say \emph{chhūṭṇo} again
\end{itemize}

\begin{tcolorbox}[breakable,colback=teal!4,colframe=teal!35!black,title={Before you move on}]
What does \mw{छूटणो} mean? (\textquotedblleft{}To depart (of a bus or a train); to get loose\textquotedblright{}.) Say it once more.
\end{tcolorbox}
\section[\texorpdfstring{\mw{पुगावणो} (pugāvṇo)}{पुगावणो (pugāvṇo)}]{\mw{पुगावणो} (pugāvṇo) — to take to a place; to see off}

\label{lesson:MW-C218-pugavno}



\begin{quote}
Before the new one: say the Marwadi for a driver, then the Marwadi for to depart.
\end{quote}

\subsection*{You'll want to know: \mw{पुगावणो}}
\textbf{\mw{पुगावणो}} — \emph{pugāvṇo} — \textquotedblleft{}to take to a place; to see off\textquotedblright{}.

\begin{quote}
\mw{हूं} \mw{काल} \mw{दूध} \mw{सहर} \mw{पुगासूं।}
\end{quote}

\begin{quote}
\emph{hū̃ kāl dūdh sahar pugāsū̃.}
\end{quote}

\begin{quote}
Tomorrow I'll take the milk to the city.
\end{quote}

Past (I …): \mw{म्हैं} \mw{दूध} \mw{सहर} \mw{पुगायो}, \emph{mhaĩ dūdh sahar pugāyo}.

\begin{cousinweb}
\textbf{\mw{पुगावणो}} (\emph{pugāvṇo}) is built from \textbf{\mw{पूगणो}}, to reach, with the causative \emph{-āv-}: to make someone or something reach a place. \textbf{\mw{पूगणो}} goes back to Sanskrit \emph{pūryate}, \textquotedblleft{}is filled\textquotedblright{}; CDIAL gives no Marwari form in this entry (CDIAL 8342).
\end{cousinweb}

\begin{grammarlens}[title={What you've built}]
One more word for saying what comes next.
\end{grammarlens}

\subsection*{Your turn}
\begin{itemize}
\raggedright
  \item \emph{Say it:} \emph{pugāvṇo}
  \item \emph{Say it:} \emph{pugāvṇo}, once more
  \item \emph{Say it:} \emph{hū̃ kāl dūdh sahar pugāsū̃.} — Tomorrow I'll take the milk to the city.
  \item \emph{Recall:} say the Marwadi for a driver, then the Marwadi for to depart, then say \emph{pugāvṇo} again
\end{itemize}

\begin{tcolorbox}[breakable,colback=teal!4,colframe=teal!35!black,title={Before you move on}]
What does \mw{पुगावणो} mean? (\textquotedblleft{}To take to a place; to see off\textquotedblright{}.) Say it once more.
\end{tcolorbox}
